\subsection{Cambio de anillo}
\label{subsection:0.1.5-fr}

\begin{env}[1.5.1]
\label{0.1.5.1-fr}
Sean $A$, $A'$ dos anillos, $\varphi$ un homomorfismo $A'\to A$, y $S$ (resp. $S'$) un subconjunto multiplicativo de $A$ (resp. $A'$), tal que $\varphi(S')\subset S$; el homomorfismo compuesto
$A'\xrightarrow{\varphi}A\longrightarrow S^{-1}A$
se factoriza en
$A'\longrightarrow S'^{-1}A'\xrightarrow{\varphi^{S'}}S^{-1}A$
en virtud de (1.2.4); se tiene $\varphi^{S'}(a'/s')=\varphi(a')/\varphi(s')$.
Si $A=\varphi(A')$ y $S=\varphi(S')$, $\varphi^{S'}$ es \emph{sobreyectivo}.
Si $A'=A$ y $\varphi$ es la identidad, $\varphi^{S'}$ no es otro que el homomorfismo $\rho_A^{S,S'}$ definido en (1.4.1).
\end{env}

\begin{env}[1.5.2]
\label{0.1.5.2-fr}
Bajo las hipótesis de (1.5.1), sea $M$ un $A$-módulo.
Existe un homomorfismo canónico funtorial
\[
  \sigma:S'^{-1}(M_{[\varphi]})\longrightarrow
  (S^{-1}M)_{[\varphi^{S'}]}
\]
de $S'^{-1}A'$-módulos, que hace corresponder a todo elemento $m/s'$ de $S'^{-1}(M_{[\varphi]})$ el elemento $m/\varphi(s')$ de $(S^{-1}M)_{[\varphi^{S'}]}$; se comprueba inmediatamente que esta definición no depende de la expresión $m/s'$ del elemento considerado.
\emph{Cuando} $S=\varphi(S')$, el homomorfismo $\sigma$ es \emph{biyectivo}.
Cuando $A'=A$ y $\varphi$ es la identidad, $\sigma$ no es otro que el homomorfismo $\rho_M^{S,S'}$ definido en (1.4.1).

Tomando en particular $M=A$, el homomorfismo $\varphi$ define en $A$ una estructura de $A'$-álgebra; $S'^{-1}(A_{[\varphi]})$ está entonces provisto de una estructura de anillo, para la cual se identifica con $(\varphi(S'))^{-1}A$, y el homomorfismo
$\sigma:S'^{-1}(A_{[\varphi]})\longrightarrow S^{-1}A$
es un homomorfismo de $S'^{-1}A'$-álgebras.
\end{env}

\begin{env}[1.5.3]
\label{0.1.5.3-fr}
Sean $M$ y $N$ dos $A$-módulos; componiendo los homomorfismos definidos en (1.3.4) y (1.5.2), se obtiene un homomorfismo
\[
  (S^{-1}M\otimes_{S^{-1}A}S^{-1}N)_{[\varphi^{S'}]}
  \longleftarrow
  S'^{-1}((M\otimes_A N)_{[\varphi]})
\]
que es un isomorfismo cuando $\varphi(S')=S$.
Igualmente, componiendo los homomorfismos (1.3.5) y (1.5.2), se obtiene un homomorfismo
\[
  S'^{-1}((\operatorname{Hom}_A(M,N))_{[\varphi]})
  \longrightarrow
  (\operatorname{Hom}_{S^{-1}A}(S^{-1}M,S^{-1}N))_{[\varphi^{S'}]}
\]
que es un isomorfismo cuando $\varphi(S')=S$ y $M$ admite una presentación finita.
\end{env}

\begin{env}[1.5.4]
\label{0.1.5.4-fr}
Consideremos ahora un $A'$-módulo $N'$ y formemos el producto tensorial $N'\otimes_{A'}A_{[\varphi]}$, que puede considerarse como un $A$-módulo poniendo
$a.(n'\otimes b)=n'\otimes(ab)$.
Existe un isomorfismo funtorial de $S^{-1}A$-módulos
\[
  \tau:(S'^{-1}N')\otimes_{S'^{-1}A'}(S^{-1}A)_{[\varphi^{S'}]}
  \xrightarrow{\sim}
  S^{-1}(N'\otimes_{A'}A_{[\varphi]})
\]
\oldpage[0\textsubscript{I}]{18}
que hace corresponder al elemento $(n'/s')\otimes(a/s)$ el elemento
$(n'\otimes a)/(\varphi(s')s)$; se comprueba por separado que, cuando se
sustituye $n'/s'$ (resp. $a/s$) por otra expresión del mismo elemento,
$(n'\otimes a)/(\varphi(s')s)$ no cambia; por otra parte, puede definirse un
homomorfismo inverso de $\tau$ haciendo corresponder a $(n'\otimes a)/s$
el elemento $(n'/1)\otimes(a/s)$: se utiliza el hecho de que
$S^{-1}(N'\otimes_{A'}A_{[\varphi]})$ es canónicamente isomorfo a
$(N'\otimes_{A'}A_{[\varphi]})\otimes_A S^{-1}A$ (1.2.5), y por tanto también a
$N'\otimes_{A'}(S^{-1}A)_{[\psi]}$, denotando por $\psi$ el homomorfismo
compuesto $a'\to\varphi(a')/1$ de $A'$ en $S^{-1}A$.
\end{env}

\begin{env}[1.5.5]
\label{0.1.5.5-fr}
Si $M'$ y $N'$ son dos $A'$-módulos, componiendo los isomorfismos (1.3.4)
y (1.5.4), se obtiene un isomorfismo
\[
  S'^{-1}M'\otimes_{S'^{-1}A'}S'^{-1}N'
  \otimes_{S'^{-1}A'}S^{-1}A
  \xrightarrow{\sim}
  S^{-1}(M'\otimes_{A'}N'\otimes_{A'}A).
\]
Igualmente, si $M'$ admite una presentación finita, se tiene, en virtud de (1.3.5) y
(1.5.4), un isomorfismo
\[
  \operatorname{Hom}_{S'^{-1}A'}(S'^{-1}M',S'^{-1}N')
  \otimes_{S'^{-1}A'}S^{-1}A
  \xrightarrow{\sim}
  S^{-1}(\operatorname{Hom}_{A'}(M',N')\otimes_{A'}A).
\]
\end{env}

\begin{env}[1.5.6]
\label{0.1.5.6-fr}
Bajo las hipótesis de (1.5.1), sea $T$ (resp. $T'$) un segundo subconjunto
multiplicativo de $A$ (resp. $A'$) tal que $S\subset T$ (resp. $S'\subset T'$)
y $\varphi(T')\subset T$. Entonces el diagrama
\[
\begin{tikzcd}[column sep=huge,row sep=large]
S'^{-1}A' \arrow[r,"\varphi^{S'}"] \arrow[d,"\rho^{T',S'}"']
  & S^{-1}A \arrow[d,"\rho^{T,S}"] \\
T'^{-1}A' \arrow[r,"\varphi^{T'}"']
  & T^{-1}A
\end{tikzcd}
\]
es conmutativo. Si $M$ es un $A$-módulo, el diagrama
\[
\begin{tikzcd}[column sep=huge,row sep=large]
S'^{-1}(M_{[\varphi]}) \arrow[r,"\sigma"] \arrow[d,"\rho^{T',S'}"']
  & (S^{-1}M)_{[\varphi^{S'}]} \arrow[d,"\rho^{T,S}"] \\
T'^{-1}(M_{[\varphi]}) \arrow[r,"\sigma"']
  & (T^{-1}M)_{[\varphi^{T'}]}
\end{tikzcd}
\]
es conmutativo. Finalmente, si $N'$ es un $A'$-módulo, el diagrama
\[
\begin{tikzcd}[column sep=huge,row sep=large]
(S'^{-1}N')\otimes_{S'^{-1}A'}(S^{-1}A)_{[\varphi^{S'}]}
  \arrow[r,"\tau","{\sim}"'] \arrow[d]
  & S^{-1}(N'\otimes_{A'}A_{[\varphi]}) \arrow[d,"\rho^{T,S}"] \\
(T'^{-1}N')\otimes_{T'^{-1}A'}(T^{-1}A)_{[\varphi^{T'}]}
  \arrow[r,"\tau"',"{\sim}"]
  & T^{-1}(N'\otimes_{A'}A_{[\varphi]})
\end{tikzcd}
\]
es conmutativo, y la flecha vertical de la izquierda se obtiene aplicando
$\rho_{N'}^{T',S'}$ a $S'^{-1}N'$ y $\rho_A^{T,S}$ a $S^{-1}A$.
\end{env}

\oldpage[0\textsubscript{I}]{19}

\begin{env}[1.5.7]
\label{0.1.5.7-fr}
Sean $A''$ un tercer anillo, $\varphi':A''\to A'$ un homomorfismo
de anillos, y $S''$ un subconjunto multiplicativo de $A''$ tal que
$\varphi'(S'')\subset S'$. Pongamos $\varphi''=\varphi\circ\varphi'$; entonces se tiene
\[
  \varphi''^{S''}=\varphi^{S'}\circ\varphi'^{S''}.
\]
Sea $M$ un $A$-módulo; se tiene evidentemente
$M_{[\varphi'']}=(M_{[\varphi]})_{[\varphi']}$; si $\sigma'$ y $\sigma''$ son
los homomorfismos definidos a partir de $\varphi'$ y $\varphi''$ como $\sigma$
se define en (1.5.2) a partir de $\varphi$, se tiene la fórmula de transitividad
\[
  \sigma''=\sigma\circ\sigma'.
\]
Finalmente, sea $N''$ un $A''$-módulo; el $A$-módulo
$N''\otimes_{A''}A_{[\varphi'']}$ se identifica canónicamente con
$(N''\otimes_{A''}A'_{[\varphi']})\otimes_{A'}A_{[\varphi]}$,
e igualmente, el $S^{-1}A$-módulo
$(S''^{-1}N'')\otimes_{S''^{-1}A''}(S^{-1}A)_{[\varphi''^{S''}]}$
se identifica canónicamente con
$((S''^{-1}N'')\otimes_{S''^{-1}A''}(S'^{-1}A')_{[\varphi'^{S''}]})\otimes_{S'^{-1}A'}(S^{-1}A)_{[\varphi^{S'}]}$.
Con estas identificaciones, si $\tau'$ y $\tau''$ son los isomorfismos definidos
a partir de $\varphi'$ y $\varphi''$ como $\tau$ se define en (1.5.4) a
partir de $\varphi$, se tiene la fórmula de transitividad
\[
  \tau''=\tau\circ(\tau'\otimes 1).
\]
\end{env}

\begin{env}[1.5.8]
\label{0.1.5.8-fr}
Sea $A$ un subanillo de un anillo $B$; para todo ideal primo
\emph{minimal} $\mathfrak p$ de $A$, existe un ideal primo minimal
$\mathfrak q$ de $B$ tal que $\mathfrak p=A\cap\mathfrak q$. En efecto,
$A_{\mathfrak p}$ es un subanillo de $B_{\mathfrak p}$ (1.3.2) y tiene
\emph{un solo} ideal primo $\mathfrak p'$ (1.2.6); como $B_{\mathfrak p}$ no se
reduce a $0$, tiene al menos un ideal primo $\mathfrak q'$ y se tiene
necesariamente $\mathfrak q'\cap A_{\mathfrak p}=\mathfrak p'$; el ideal primo
$\mathfrak q_1$ de $B$, preimagen de $\mathfrak q'$, satisface por tanto
$\mathfrak q_1\cap A=\mathfrak p$, y \emph{con mayor razón} se tiene
$\mathfrak q\cap A=\mathfrak p$ para todo ideal primo minimal $\mathfrak q$
de $B$ contenido en $\mathfrak q_1$.
\end{env}
